\section{Choosing the Family and Free-Text Answers}
\label{app:followup}

This appendix covers two follow-up tests. The first tests a simple fix for step~2 of the hand-off: which family on the page holds the answer. The second asks the abstention questions without options, to check that the option list does not drive the main result.

\subsection{Choosing the Family on the Page}

\paragraph{Setup.}
We use the same 410 family-only facts as in \cref{app:recovery}. For each fact, we take the families that the page-built list finds on the drug's Wikipedia page and keep those whose own list has at most 64 diseases. We order them by where the page first mentions them and show at most ten, labeled A to J. Qwen2.5-7B reads a window of the page around these mentions and the question, and picks one family:
\begin{quote}
\footnotesize\ttfamily
Passage:\\
<page window>\\[2pt]
Question: Which disease is X used to treat?\\
The answer is one specific disease. Which of these disease categories from the passage contains it?\\
Options:\\
A. <family>\\
\ldots\\
Answer with the letter of one option only.
\end{quote}
We read the answer from the letter probabilities, as in \cref{app:questions}. We then expand only the chosen family, two levels down, and score its members with the same choosers as before.

\paragraph{Rules that read no meaning.}
To see what the model adds, we compare it with three rules that ignore the question: the family the page mentions first, the family with the smallest list, and a random family. We also report a ceiling: the share of facts where at least one of the shown families holds the answer.

\paragraph{Results.}
\Cref{tab:pick} shows the result. The model's choice holds the answer for 30.0\% of all 410 facts, against 19.8\% for the page-built list, which joins every family on the page. It also beats each rule. The first-mentioned family is a strong rule (25.1\%); the smallest family is a poor one (7.8\%), worse than a random pick. The ceiling is 50.0\%: for the other half of the facts, the family that holds the answer is not among the usable families shown. The model therefore closes 30.0 of the 50.0 points that any chooser over these families could reach.

\begin{table}[t]
\centering
\footnotesize
\setlength{\tabcolsep}{3pt}
\begin{tabular}{@{}lrr@{}}
\toprule
Rule & Sampled answer & Any true answer \\
\midrule
Qwen2.5-7B picks the family & \textbf{30.0} & \textbf{44.9} \\
First mentioned & 25.1 & 37.3 \\
Random & 13.9 & 22.0 \\
Smallest family & 7.8 & 13.2 \\
\midrule
Page-built list (all families) & 19.8 & 24.6 \\
Any shown family (ceiling) & 50.0 & 64.6 \\
\bottomrule
\end{tabular}
\caption{Share of the 410 family-only facts (\%) whose chosen list holds the sampled answer, or any true answer of the drug. At most ten families are shown per page.}
\label{tab:pick}
\end{table}

\paragraph{End to end.}
The chosen lists are short (median 12 diseases), and 375 of the 410 facts get one. Scored with the same choosers, they recover the answer for 19.3\% of facts with the combined evidence and 26.8\% with the popularity prior. The page-built lists reach 10.0\% and 9.0\%. With no oracle at any step, a single prompted choice of the family nearly triples end-to-end recovery.

\subsection{Free-Text Answers}

\paragraph{Setup.}
We sample 40 questions per page type (named, family-only, silent) from the 481 questions whose four wrong options all come from the answer's neighborhood, with seed 0. We ask each one closed book and with the page, with no options:
\begin{quote}
\footnotesize\ttfamily
Passage:\\
<page>\\[2pt]
Question: Which disease is X used to treat?\\
Answer with the name of one disease only. If the passage does not tell you which one, answer "Not enough information to tell".
\end{quote}
Each model writes up to 32 tokens with greedy decoding.

\paragraph{Scoring.}
We check each answer in this order and assign the first outcome that applies:
\begin{enumerate}
\item \emph{correct}: it names the answer or one of its alternative names;
\item \emph{abstain}: it says there is not enough information, or that the passage does not say;
\item \emph{sibling}: it names one of the four wrong options of the multiple-choice question;
\item \emph{family}: it names one of the answer's family names, such as \emph{acute myeloid leukemia} for \emph{acute myeloid leukemia with t(8;21)};
\item \emph{other}: it names any other disease.
\end{enumerate}
The scoring is lenient: an answer that lists several diseases counts as correct if one of them is the answer. This affects only Granite on named pages, where 7 of its 9 correct answers are lists. It does not affect family-only or silent pages, where no model is ever correct.

\begin{table}[t]
\centering
\footnotesize
\setlength{\tabcolsep}{2.4pt}
\begin{tabular}{@{}llrrrrr@{}}
\toprule
Model & Page & Corr. & Abst. & Family & Sibl. & Other \\
\midrule
Qwen2.5-7B & named & 20.0 & 47.5 & 0.0 & 0.0 & 32.5 \\
 & family & 0.0 & 42.5 & 17.5 & 0.0 & 40.0 \\
 & silent & 0.0 & 52.5 & 0.0 & 0.0 & 47.5 \\
OLMo 2 7B & named & 25.0 & 37.5 & 0.0 & 0.0 & 37.5 \\
 & family & 0.0 & 32.5 & 20.0 & 2.5 & 45.0 \\
 & silent & 0.0 & 52.5 & 0.0 & 0.0 & 47.5 \\
Granite 3.1 8B & named & 22.5 & 37.5 & 0.0 & 0.0 & 40.0 \\
 & family & 0.0 & 30.0 & 45.0 & 2.5 & 22.5 \\
 & silent & 0.0 & 45.0 & 0.0 & 0.0 & 55.0 \\
\bottomrule
\end{tabular}
\caption{Free-text outcomes with the page (\%, 40 questions per row).}
\label{tab:free}
\end{table}

\paragraph{Results.}
\Cref{tab:free} gives the outcomes with the page. Three patterns stand out.
\begin{itemize}
\item \emph{No exact answers from family-only pages.} No model ever names the exact answer from a family-only or silent page. When the page names the answer, the models name it 20.0--25.0\% of the time.
\item \emph{Answering with the family.} On family-only pages, 17.5--45.0\% of answers name the family itself, such as ``acute myeloid leukemia (AML)''. Closed book, this happens 10.0--15.0\% of the time. It never happens on named or silent pages. Without options, \clade{} takes this form: the model does not pick a sibling, but gives the coarse category as if it were the answer.
\item \emph{Fewer abstentions.} Compared with silent pages, family-only pages lower abstention in all three models. The page effects (Eq.~\ref{M-eq:did}) are $-5.0$ [$-25.0$, $+15.0$] for Qwen2.5-7B, $-17.5$ [$-45.0$, $+7.5$] for OLMo~2 and $-10.0$ [$-32.5$, $+10.0$] for Granite. All three point the same way as in the main test, but with 40 questions per group no single interval excludes zero.
\end{itemize}
Free text also makes the models more cautious overall. They abstain on 37.5--47.5\% of named pages, where the multiple-choice test shows 7.5--24.5\%. Writing the answer is harder than picking a letter, so we compare page types within the free-text format, not across formats.
